\chapter{Modelo predictivo}\label{cap:prediccion}

El modelo de efectos se construyó para interpretar. Si el objetivo es predecir la mortalidad de
condados nuevos, el criterio cambia: lo que importa es el error fuera de la muestra, y pueden
convenir modelos con más variables, regularizados o no lineales \parencite{hastie2009}. Este
capítulo compara ocho modelos con un protocolo que evita estimaciones optimistas del error.

\section{Protocolo}

\begin{enumerate}
  \item Se reserva una partición de prueba de \res{n-prueba} condados (20\,\%, estratificada por
  región) que no se usa para ajustar ni para elegir. Los \res{n-entrenamiento} restantes forman el
  entrenamiento.
  \item En entrenamiento se comparan los modelos por validación cruzada de 10 pliegues. Los
  hiperparámetros (penalización de ridge, lasso y red elástica; profundidad y tasa de aprendizaje de
  los árboles) se eligen dentro de cada pliegue con una segunda validación cruzada (validación
  cruzada \emph{anidada}), y la imputación de ausentes y la tipificación se ajustan también dentro
  del pliegue, para no filtrar información de la parte de validación.
  \item Se repite con validación cruzada agrupada por estado: cada pliegue deja fuera estados
  completos. Mide la capacidad de predecir condados de estados no vistos; una partición aleatoria es
  optimista porque los condados de un mismo estado se parecen.
  \item El modelo con menor RMSE en validación cruzada se evalúa una única vez en la prueba.
\end{enumerate}

Los modelos comparados son: la media del entrenamiento (referencia), MCO con la especificación
del modelo de efectos, MCO con todas las candidatas, ridge, lasso y red elástica (que pueden usar
también indicadoras de estado), un bosque aleatorio y una potenciación del gradiente sobre
histogramas.

\section{Resultados}

\begin{table}[htbp]
\centering
\footnotesize
\caption{Comparación de modelos predictivos: validación cruzada en entrenamiento y evaluación en prueba}\label{tab:prediccion}
\begin{tabularx}{\linewidth}{X r r r r r r r}
\toprule
Modelo & \multicolumn{3}{c}{VC aleatoria (10)} & VC estado & \multicolumn{3}{c}{Prueba} \\ \cmidrule(lr){2-4}\cmidrule(lr){5-5}\cmidrule(lr){6-8} & RMSE & D.\,T. & $R^2$ & RMSE & RMSE & $R^2$ & MAE \\
\midrule
Referencia (media) & \num{27.90} & \num{1.11} & \num{-0.003} & \num{27.93} & \num{27.10} & \num{-0.002} & \num{21.16} \\
MCO con la especificación del modelo de efectos & \num{19.11} & \num{1.44} & \num{0.527} & \num{19.43} & \num{18.18} & \num{0.549} & \num{13.85} \\
MCO con todas las variables & \num{18.61} & \num{1.37} & \num{0.551} & \num{20.36} & \num{17.44} & \num{0.585} & \num{13.26} \\
Ridge$^\star$ & \num{18.59} & \num{1.37} & \num{0.553} & \num{19.68} & \num{17.36} & \num{0.589} & \num{13.18} \\
Lasso & \num{18.60} & \num{1.36} & \num{0.552} & \num{19.65} & \num{17.35} & \num{0.589} & \num{13.17} \\
Red elástica & \num{18.60} & \num{1.36} & \num{0.552} & \num{19.65} & \num{17.36} & \num{0.589} & \num{13.19} \\
Bosque aleatorio & \num{18.95} & \num{1.01} & \num{0.536} & \num{19.93} & \num{17.60} & \num{0.577} & \num{13.22} \\
Potenciación del gradiente (histogramas) & \num{19.03} & \num{0.90} & \num{0.532} & \num{19.81} & \num{17.61} & \num{0.577} & \num{13.36} \\
\bottomrule
\end{tabularx}
\par\smallskip{\footnotesize\raggedright $^\star$: modelo elegido por la regla configurada. La prueba se usa una sola vez, tras elegir.\par}
\end{table}


\figura{validacion-cruzada}{RMSE de cada modelo en validación cruzada aleatoria (media y
desviación típica entre pliegues) y agrupada por estado.}

El mejor modelo en validación cruzada es \res{modelo-elegido} (RMSE \res{cv-rmse-ridge}); lasso y
red elástica quedan prácticamente empatados (\res{cv-rmse-lasso}), y la regla de un error típico,
que elegiría el modelo más sencillo cuyo error no supere al mínimo en más de una desviación típica,
elegiría \res{modelo-1ee}. Los modelos de árboles no mejoran a los lineales regularizados: la
relación es esencialmente lineal y aditiva, como sugería el diagnóstico. En la partición de prueba
el modelo elegido obtiene un RMSE de \res{test-rmse-elegido} \res{test-rmse-ic} (intervalo
bootstrap), $R^2=\res{test-r2-elegido}$ y un error absoluto medio de \res{test-mae-elegido}; su
ventaja sobre el MCO con la especificación de efectos es pequeña pero significativa (diferencia de
RMSE \res{test-dif-ic}), atribuible a las covariables adicionales y, sobre todo, a las
indicadoras de estado, que el modelo de efectos no usa. Ambos reducen en torno a un tercio el error
de la predicción ingenua por la media (\res{test-rmse-referencia}).

\figura{prueba}{Partición de prueba. Izquierda: mortalidad observada frente a la predicha por el
modelo elegido, por región. Derecha: cobertura de los intervalos conformales estándar y
normalizados según el tamaño del condado.}

Dos matices sobre la honestidad de estas cifras. Primero, el RMSE de prueba es algo menor que el de
validación cruzada por azar de la partición: el intervalo bootstrap lo cuantifica. Segundo, la
especificación del modelo de efectos se seleccionó con todos los condados, de modo que su error en
prueba es ligeramente optimista; el resto de modelos se ajustó sólo con el entrenamiento.

Cuando hay que predecir estados no vistos (validación agrupada), el error aumenta para todos los
modelos y el mejor pasa a ser el \res{modelo-mejor-gcv} (RMSE \res{gcv-rmse-mco-efectos}): las
indicadoras de estado no sirven para estados nuevos, y un modelo parsimonioso generaliza mejor.

\section{Importancia de las variables}

La importancia por permutación mide cuánto empeora el RMSE en prueba al desordenar cada variable.
Las tres más importantes para el modelo elegido son \res{importancia1}, \res{importancia2} y
\res{importancia3}, en coherencia con el modelo de efectos.

\figura{importancia}{Importancia por permutación en la partición de prueba (aumento medio del RMSE
y su desviación típica en 15 permutaciones).}

\section{Intervalos de predicción conformales}

Para acompañar cada predicción con un intervalo sin suponer normalidad se usó la predicción
conformal \parencite{lei2018}: con los residuos fuera de pliegue del entrenamiento se calcula el
cuantil que garantiza, para un condado nuevo intercambiable con los de entrenamiento, una cobertura
de al menos el \res{conf-nominal}\,\%. El intervalo estándar tiene la misma anchura
(\res{conf-anchura-std}) para todos los condados y cubre el \res{conf-std}\,\% en prueba, pero de
forma desigual: \res{conf-std-peq}\,\% en los condados pequeños y \res{conf-std-gra}\,\% en los
grandes. El intervalo \emph{normalizado} divide los residuos por $\hat\sigma(n)$, la desviación
típica que da la función de varianza, y se ensancha donde la tasa es más ruidosa: cubre el
\res{conf-norm}\,\% globalmente, con una cobertura más equilibrada (\res{conf-norm-peq}\,\% en los
pequeños, \res{conf-norm-gra}\,\% en los grandes) y anchuras de \res{conf-anchura-norm-peq} a
\res{conf-anchura-norm-gra} (\cref{tab:conformal}).

\begin{table}[htbp]
\centering
\small
\caption{Cobertura empírica en prueba de los intervalos de predicción conformales al \num{90}\,\%}\label{tab:conformal}
\begin{tabular}{l r r r r}
\toprule
Condados & \multicolumn{2}{c}{Estándar} & \multicolumn{2}{c}{Normalizado} \\ \cmidrule(lr){2-3}\cmidrule(lr){4-5} & Cobertura (\%) & Anchura & Cobertura (\%) & Anchura \\
\midrule
Pequeños & \num{81.9} & \num{60.0} & \num{88.2} & \num{71.4} \\
Medianos & \num{95.6} & \num{60.0} & \num{91.6} & \num{51.1} \\
Grandes & \num{96.1} & \num{60.0} & \num{92.6} & \num{44.5} \\
\midrule
Total & \num{91.1} & \num{60.0} & \num{90.8} & \num{55.7} \\
\bottomrule
\end{tabular}

\end{table}

